% ---
% title: Difference with KL in Loss and reward
% date: 2026/2/9
% tag: RL
% ---
\definecolor{mdc1}{HTML}{D32F2F}
\definecolor{mdc2}{HTML}{FFFFFF}
\definecolor{mdc3}{HTML}{0000FF}
\definecolor{mdc4}{HTML}{00AA00}
\definecolor{mdc5}{HTML}{FF0000}
为什么在GRPO中进行微调时，在损失（Loss）里添加散度项（Reverse KL Divergence)本质上和在奖励（Reward）中添加散度项是\colorbox{mdc1}{不同的}，并且会造成Bias？

\hrule

\section{一、优化目标}

\[
J(\theta) = \mathbb{E}_{y\sim\pi_\theta} \big[ \widehat{\mathrm{KL}}(y) \big]
\]

这是一个\colorbox{mdc1}{\textcolor{mdc2}{参数在分布里}}的期望。

\section{二、关键}

\begin{quote}
\textbf{当参数} \(\theta\) \textbf{同时出现在：}
\end{quote}

\begin{itemize}
  \item \begin{quote}
  分布里（\(\pi_\theta\)）
  \end{quote}
  \item \begin{quote}
  被积函数里（\(\widehat{\mathrm{KL}}(y)\)）
  \end{quote}
\end{itemize}

\begin{quote}
\textbf{对期望求导一定会产生两项}
\end{quote}

这是一个数学事实，不是实现细节。

\section{三、一步一步来：这两项是怎么来的？}

把期望写成最原始的形式：

\[
J(\theta) = \sum_y \pi_\theta(y)\; \widehat{\mathrm{KL}}(y)
\]

现在对 \(\theta\) 求导：

\[
\nabla_\theta J = \nabla_\theta \sum_y \pi_\theta(y)\; \widehat{\mathrm{KL}}(y)
\]

\subsection{用乘法求导法则（这是关键）}

\[
\nabla_\theta(\pi_\theta \cdot \widehat{\mathrm{KL}}) = (\nabla_\theta \pi_\theta)\widehat{\mathrm{KL}} + \pi_\theta (\nabla_\theta \widehat{\mathrm{KL}})
\]

所以：

\[
\nabla_\theta J = \sum_y \Big[ \underbrace{\pi_\theta(y)\nabla_\theta \widehat{\mathrm{KL}}(y)}_{\text{第一项}} + \underbrace{\widehat{\mathrm{KL}}(y)\nabla_\theta \pi_\theta(y)}_{\text{第二项}} \Big]
\]

\section{四、化简}

用恒等式：

\[
\nabla_\theta \pi_\theta(y) = \pi_\theta(y)\nabla_\theta \log \pi_\theta(y)
\]

代回去：

\[
\nabla_\theta J = \mathbb{E}_{y\sim\pi_\theta} \big[ \nabla_\theta \widehat{\mathrm{KL}}(y) + \widehat{\mathrm{KL}}(y)\nabla_\theta \log \pi_\theta(y) \big]
\]

\section{五、自动微分能算哪一项？}

\subsection{KL 当 loss（GRPO / Loss-based KL）}

代码逻辑是：

\begin{verbatim}
y = sample(pi_theta)      # 采样
kl_hat = KL_hat(y)        # 计算 KL 估计
loss = kl_hat
loss.backward()
\end{verbatim}

\subsection{自动微分看到的计算图是：}

\[
θ → log \pi_\theta (y) → \hat{KL} → loss
\]

\colorbox{mdc3}{\textcolor{mdc2}{注意：y 是采样出来的整数 token}}

\begin{itemize}
  \item \(y\) 不是 \(\theta\) 的函数
  \item \(y\) 在计算图里是常数（detached）
\end{itemize}

\subsection{自动微分只能算到什么？}

自动微分只能看到：\(\text{loss} = \widehat{\mathrm{KL}}(y)\)，所以它只能算：\(\nabla_\theta \widehat{\mathrm{KL}}(y)\)。

\colorbox{mdc4}{\textcolor{mdc2}{然后对样本做平均，求出}}：\(\mathbb{E}_{y\sim\pi_\theta} \big[ \nabla_\theta \widehat{\mathrm{KL}}(y) \big]\) ，\colorbox{mdc5}{\textcolor{mdc2}{但是无法得到}}：\(\mathbb{E}_{y\sim\pi_\theta} \big[ \widehat{\mathrm{KL}}(y) \nabla_\theta \log \pi_\theta(y) \big]\)。

👉 来自 \(\pi_\theta(y)\) \textbf{对} \(\theta\) \textbf{的导数}

但在代码里：\(y = sample(\pi_{\theta})\)

\begin{itemize}
  \item sampling \textbf{不是可微算子}
  \item 自动微分 \textbf{不知道 y 的概率密度}
  \item 所以它 \textbf{完全不知道}
  
  \begin{quote}
  “改变 \(\theta\) 会改变 \(y\) 被采样到的概率”
  \end{quote}
\end{itemize}

📌 \textbf{这条梯度路径在计算图里根本不存在}

\begin{quote}
\(\nabla_\theta \mathbb{E}_{y\sim\pi_\theta}[\cdot]\)与\(\mathbb{E}_{y\sim\pi_\theta}[\nabla_\theta \cdot]\)而这两者 \textbf{只有在分布不依赖参数}时才相等，这里分布恰恰依赖 \(\theta\)。
\end{quote}

\section{六、为什么 RL 不会犯这个错？}

因为 RL \textbf{不是靠自动微分穿过采样}，而是\textbf{手动写出了那一项}：

\[
\nabla_\theta \mathbb{E}_{y\sim\pi_\theta}[R(y)] = \mathbb{E}_{y\sim\pi_\theta} \big[ R(y)\nabla_\theta \log \pi_\theta(y) \big]
\]

当你把：\(R(y) = -\beta \widehat{\mathrm{KL}}(y)\)

那缺失的一项 \textbf{被显式加回来了}。

\section{七、总结}

\begin{quote}
👉 自动微分只看到了\(\nabla_\theta \widehat{\mathrm{KL}}\)

👉 却完全没看到\(\nabla_\theta \pi_\theta\)

因此 \(\text{Loss-based KL} ⇒ \text{biased gradient}\)
\end{quote}
